\chapter{Theoretical phenomenon and model}
\label{phenomenon}

\section{System description}
\begin{figure}[!htb]
  \centering
  \subfloat[System schematic\label{Schematic-subfig-1}]{
    \includegraphics[totalheight=3cm,trim={0 0 0 0cm},clip]{parallel-wall-droplet.png}
  }
  \\
  \subfloat[Quadrupolar flow-field and particle interactions\label{Schematic-subfig-2}]{
    \includegraphics[totalheight=4cm,trim={0 3cm 0 3cm},clip]{quadrupole_whole.png}
  }
  \\
  \subfloat[Swapping trajectory effect and particle interactions\label{Schematic-subfig-3}]{
    \includegraphics[totalheight=6cm,trim={0 2cm 2cm 2cm},clip]{swapping_deform_whole.png}
  }
  
  \caption{System description and mechanisms for particle interactions - Subfig.~\ref{Schematic-subfig-1} shows the system schematic where a single droplet layer suspended between two parallel walls, the mid-plane is at z=0 and flow-direction is taken as x axis. Subfig.~\ref{Schematic-subfig-2} and ~\ref{Schematic-subfig-3} describe the streamlines (from FD/FT simulations) around a confined drop in x-y and x-z plane respectively, when subjected to an incident couette flow and their effect on particle interactions.}
  \label{Interaction Schematic}    
  
\end{figure}

It was observed that in a system of particles strongly confined between two parallel walls, their interaction is driven by mainly three types of hydrodynamic mechanisms-a) Hele--Shaw quadrupoles, b)swapping-trajectory effect, and c) deformation-induced particle drift towards the middle of the flow-gradient(xz) plane. The Hele--Shaw quadrupoles which are formed by the scattering of incident flow from the drops cause alignment of neighbouring drops and subsequent attraction between them. The Swapping-trajectory effect drives drops out of the mid-plane in the gradient-direction and causes an effective repulsion between them thereby maintaining a fixed inter-drop separation and finally the deformation-induced drift imparts a restoring velocity to the out-of-plane deformed particles opposite in direction to the swapping-trajectory effect i.e. directed towards the flow-gradient mid-plane.

Under strong confinement, a flow-field can be divided into a combination of incident($\mathbf v_0$) and scattered flow($\delta \mathbf{v}_{1}$). The latter comprises the velocities that give rise to the above mentioned inter-particle interactions. Hence, the total flow field can be written as,
\begin{equation}
  \mathbf{v}_1 = \mathbf{v}_0 + \delta \mathbf{v}_{1}
  \label{total flow-field}
\end{equation}
In the above equation,  $\mathbf v_0 = \dot{\gamma}z \uvec{e}_x$, where $\dot{\gamma}z$ is the shear rate, $\uvec{e}$ is the unit vector in the flow-direction and z in the gradient direction coordinate which have been schematically displayed in subfig.~\ref{Schematic-subfig-1}. It is due to the second term on the right hand side of Eq.~\ref{total flow-field} that the interparticle interactions and their respective relative motion take palce.

The left frame in Fig.~\ref{Schematic-subfig-2} displays the flow field~\ref{total flow-field} for an individual drop suspended in the mid-plane of the channel under the influence of an incident couette flow. It can be observed that, with the drop at the centre, the flow field is symmetric in all four quadrants just like a Hele-Shaw quadrupole. This supports the fact that in flows between highly confined parallel walls, the scattered flow-field generated by the particles has a far-field 2-D form at sufficiently larger distances. Additionally, for shear flow, the 2-D quadrupoles are the leading order contribution in the multipolar solutions. The right frame in Fig.~\ref{Schematic-subfig-2} shows the role played by the quadrupolar contributions in interparticle  interactions that promote drop alignment and attraction thereby driving the chain formation.

The top frame in Fig.~\ref{Schematic-subfig-3} shows the flow field of the same particle mentioned above in the flow-gradient plane. It can be observed that closer to the channel mid-plane, there are oppositely directed reversed streamlines on either side of the drop which give rise to the swapping-trajectory effect. As shown in the bottom frame of Fig.~\ref{Schematic-subfig-3}, this phenomenon drives particles out of the mid-plane which are then carried away by the streamlines in opposite directions thereby causing an effective repulsion between them. This plays an important role in determining the interparticle spacings in a chain structure. In deformable particles, a deformation-induced restoring effect imparts velocities directed towards the mid-plane. This effect dominates over the swapping-trajectory effect as drop separation increases beyond a certain distance(say $h_{z}^{bal}$).

The above mentioned hydrodynamic interactions have been modelled using a simple model which is discussed in the following sections and their effects on the the evolution of a binary system followed by that of a linear multi-particle system are described in Chapter.~\ref{Results}.

\section{Modelling of quadrupolar interactions - or flow-vorticity plane dynamics}

In far-field regime, flow scattered from a particle in a confined parallel-wall geometry tends to have a quasi-two-dimensional Hele--Shaw form given by
\begin{equation}
  \delta \mathbf{v}_{1} = -\frac{1}{8}(H-2z)(H+2z) \nabla p_1(\bm{\rho})
\end{equation}
where $\eta$ is the fluid viscosity, H is the channel height, and $\bm\rho = \textit{x} \uvec{e}_x + \bm\rho = \textit{y}\uvec{e}_y$ is the position in the flow-vorticity plane with respect to the particle center.
In the case of couette flow, since the system is symmetric when approached from either side in the flow direction, the leading order multipolar contribution is the Hele--Shaw quadrupoles given by
\begin{equation}
  \delta p_1(\bm\rho) = \dot{\gamma}\eta q p_q(\bm\rho)
\end{equation}
where,
\begin{equation}
  p_q(\bm\rho) = -\frac{d^2(x^2 - y^2)}{4 \rho^4}
\end{equation}
is the non-dimensional quadrupolar pressure field, q is the quadrupolar moment, $\rho = \left|\bm{\rho}\right|$ and d is the drop diameter. Thus the velocity of a particle in the system after incorporating the pairwise superposition approximation is given by
\begin{equation}
  \mathbf{U}_{qi} = \dot{\gamma}dQ \sum_{i \neq j}\mathbf{v}_q(\bm{\rho}_{ij})
\end{equation}
where,
\begin{equation}
  \mathbf{v}_q(\bm\rho) = d^3 \nabla \frac{x^2 - y^2}{4 \rho^4}
\end{equation}
is the normalized Hele--Shaw quadrupolar velocity flow field and $\bm{\rho}_{ij} = \bm{\rho}_{i}-\bm{\rho}_{j}$ is the relative position of a particle pair and $Q = \mu q {H/d}^2$ is the dimensionless quadrupolar moment.

Additionally, since an unbalanced attractive force due to quadrupolar interactions alone can cause particle overlap, hence we model the excluded-volume interactions using a simple repulsion potential
\begin{equation}
  \mathbf{u}_R(\bm\rho) = -\nabla V(\rho) 
\end{equation}
where
\begin{equation}
 V(\rho)  = A_0 \rho^{-m}e^{-\kappa_0\rho}
\end{equation}
is a general form of the Yukawa potential(m=8 and $\kappa_0d=6$). The amplitude $A_0$ has a value such that the quadrupolar and excluded volume attactive and repulsive forces respectively, balance each other at particle contact.

\section{Modelling of swapping trajectory interactions}

The swapping trajectory effect is caused between a particle pair when the incident flow scattered from one particle gets reflected from the confining wall towards the other particle\cite{Reference27}. This reflected flow drives the particle transversely and moves them away from each other via oppositely directed streamlines of the couette flow. This interaction is reflected in the reversing streamlines around a single confined particle and also the particle trajectories as shown in Fig.~\ref{Schematic-subfig-3}.

The traverse velocity($u_{iz}(\bm{\rho}_{ij},z_i)$) of the deformable drops are comprised of two competing phenomena namely(also described in Section.3.1) the swapping-trajectory velocity and the deformation induced individual particle drift velocity towards flow-gradient mid-plane(also known as the self-centering effect). It can therefore be written as
\begin{equation}
 U_{iz}(\bm{\rho}_{ij},z_i) = u_{iz}^{sw}(\bm{\rho}_{ij}) + u_{iz}^c(z_i).
\end{equation}
Due to the observation that the displacements of the particles from the mid-plane($z=0$) is small, it is assumed that the deformation induced restoring velocity has the form
\begin{equation}
  u_{iz}^c(z_i) = -\beta \dot{\gamma}z_i.
  \label{deformation velocity}
\end{equation}
where $\beta$ depends on particle deformation and confinement of the system($H/d$). $\beta$ was calculated from the FD/FT simulation of a single drop placed at a finite distance above the mid-plane, in a periodic unit cell. Details of which are discussed in the Appendix.

Considering the far-field form of the wall reflected flow responsible for the cross-channel migration near a wall\cite{Reference27} and based on the fact that the flow fields in a parallel wall system decay exponentially(with the exception of Hele--Shaw fields) a simple approximate form for the swapping trajectory velocity is proposed
\begin{equation}
  \mathbf{u}_{ix}^{sw}(\bm{\rho}_{ij}) = \dot{\gamma}dAu_0^{sw}(\bm{\rho}_{ij}) \uvec{e}_x
  \label{transverse swapping}
\end{equation}
where
\begin{equation}
  u_0^{sw}(\bm{\rho}) = \frac{d^3x \bar{z}}{(rev_h + e^{\kappa(\rho-\rho_0)})(b^2 +  \left|\bm{\rho}\right|^2)^2}
  \label{swapping form}
\end{equation}
where A, $\kappa$, $rev_h$, b and $\rho_0$(distance of transition from polynomial to exponential decay of the flow),are parameters sued to match a data from the FD/FT simulations and $\bar{z}$ is the transverse location of the particle's image.
%Since the axial repulsive component of the swapping-trajectory effect is due to the convection of a particle due to its non-zero transverse position in the incident couette flow, hence for a given axial-position($z_i$) we can represent it as
% \begin{equation}
%  \mathbf{u}_{xi}^{sw}(z_i) = \alpha \dot{\gamma}z_i \uvec{e}_x
% \end{equation}
% where $\alpha$ is the mobility coefficient. Taking this fact into consideration with Eq.~\ref{swapping form} the form o the swapping-trajectory repulsion is taken as
In the flow-direction, 
\begin{equation}
 \mathbf{u}_{xi}^{sw}(z_i) = \dot{\gamma}dBu_0^{sw}(\bm{\rho}_{ij}) \uvec{e}_x.
\end{equation}
where B=$f(A \alpha)$.

The above mentioned parameters in Eqs.~\ref{transverse swapping}-~\ref{deformation velocity} have been obtained for an interacting pair of partcles. In a multibody system, the swapping-trajectory velocities are obtained by the superposition of pairwise results.
Therefore, the final expression for the hydrodynamic interactions in the flow direction is represented as
\begin{equation}
 \mathbf{U}_{xi} = \dot{\gamma}d \sum_{i \neq j}[Q\mathbf{v}_q(\bm{\rho}_{ij}) + Bu_0^{sw}(\bm{\rho}_ij) \uvec{e}_x + \mathbf{u}_R(\bm{\rho}_{ij})].
  \label{total flow}
\end{equation}

\section{System parameters for quadrupolar and swapping-trajectory interactions}
A calculation of the Q and the swapping trajectory parameters was performed to match the benchmark binary particle system dynamics in a periodic unit cell which were linearly arranged along the flow direction. While the quadrupolar moment, Q was evaluated using the far-field data of the FD/FT simulations, the swapping-trajectory parameters were obtained from the near field dynamics. Details of the mathcing are in the Appendix and the values of the parameters used in our model are listed in the Table.~\ref{tab:table1}

\begin{table}[htb!]
  \caption{Parameters obtained for the 3-D model through benchmark simulations}
  \centering 
  \label{tab:table1}
  \begin{tabular}{c c c c c c} % <-- Alignments: 1st column left, 2nd middle and 3rd right, with vertical lines in between
  \hline\hline                        %
  \textbf{$Q$} & \textbf{$\kappa d$} & \textbf{$b/d$} & \textbf{$\rho_0/d$} & \textbf{$rev_h$} & \textbf{$B$} \\
  % $Q$ & $\kappa d$ & $b/d$ & $\rho_0$ & B & $rev_h$\\
  \hline
  0.0816 & 9 & 1 & 1.90 & 0.5 & 2.26\\
  \hline
\end{tabular}
\end{table}